\documentclass[11pt]{article}
\usepackage[margin=1in]{geometry}
\usepackage{longtable}
\usepackage{booktabs}
\usepackage{array}
\usepackage{hyperref}
\usepackage{enumitem}
\usepackage{caption}

\title{Implementation Audit: Sparse Signal-Field SAE Plan vs Codebase}
\author{Project Audit Report}
\date{June 1, 2026}

\begin{document}
\maketitle

\section{Purpose and Scope}
This document provides a thorough, evidence-based audit of the project state against the high-level 8-experiment SAE plan.
It addresses the four requested questions:
\begin{enumerate}[leftmargin=*]
    \item Which of the 8 planned experiments have been implemented, and which have not.
    \item A thorough summary of implementation changes from the high-level plan to the actual code.
    \item A thorough summary of experiments performed after those implementation changes.
    \item Results.
\end{enumerate}

Audit date: \textbf{June 1, 2026}.\newline
Primary project root: \texttt{Vision\_Sparsity/CNN-SAE}.

\section{Primary Evidence Reviewed}
\subsection{High-level planning artifacts}
\begin{itemize}[leftmargin=*]
    \item \texttt{high-level/sparse\_signal\_sae\_experiment\_plan\_high\_accuracy\_conv.tex}
    \item \texttt{high-level/sparse\_signal\_sae\_experiment\_plan\_high\_accuracy\_conv.pdf}
\end{itemize}

\subsection{Project guidance and status artifacts}
\begin{itemize}[leftmargin=*]
    \item \texttt{README.md}
    \item \texttt{TODO.md}
    \item \texttt{docs/PHASE1\_RESULTS.md}
    \item \texttt{docs/PHASE2\_RESULTS.md}
    \item \texttt{docs/PHASE3\_RESULTS.md}
    \item \texttt{docs/PHASE3b\_RESULTS.md}
    \item \texttt{report/REPORT.md}
\end{itemize}

\subsection{Execution and result artifacts}
\begin{itemize}[leftmargin=*]
    \item \texttt{scripts/run\_grid.py}, \texttt{scripts/run\_phase3.sh}, tier scripts
    \item \texttt{runs/phase1}, \texttt{runs/phase2}, \texttt{runs/phase3}, \texttt{runs/phase3b}, \texttt{runs/phase3b\_bptt}
    \item \texttt{runs/intervene\_result.json}, \texttt{runs/audit/audit\_Q3.json}, \texttt{runs/audit/audit\_Q5.json}
    \item \texttt{logs/baselines\_pca.json}, backbone done files, and tier logs
\end{itemize}

\section{The Planned 8 Experiments (from the high-level spec)}
The high-level spec defines an 8-cell factorial:
\begin{itemize}[leftmargin=*]
    \item 2 training protocols: Independent (Family I) vs Incremental dependent (Family II)
    \item 2 SAE types: Field vs Vector
    \item 2 mask types: Global vs Receptive-field-local (RF-local)
\end{itemize}

Mapped IDs in the spec:
\begin{itemize}[leftmargin=*]
    \item E1: Independent + Field + Global
    \item E2: Independent + Field + RF-local
    \item E3: Independent + Vector + Global
    \item E4: Independent + Vector + RF-local
    \item E5: Incremental + Field + Global
    \item E6: Incremental + Field + RF-local
    \item E7: Incremental + Vector + Global
    \item E8: Incremental + Vector + RF-local
\end{itemize}

\section{1) Implementation Status of E1-E8}
\subsection{Direct artifact check}
Directory evidence in \texttt{runs/}:
\begin{itemize}[leftmargin=*]
    \item E1 directories found: 10
    \item E2 directories found: 10
    \item E3 directories found: 15
    \item E4 directories found: 10
    \item E5 directories found: 0
    \item E6 directories found: 0
    \item E7 directories found: 0
    \item E8 directories found: 0
\end{itemize}

\subsection{Conclusion}
\begin{longtable}{>{\raggedright\arraybackslash}p{0.08\linewidth} >{\raggedright\arraybackslash}p{0.25\linewidth} >{\raggedright\arraybackslash}p{0.18\linewidth} >{\raggedright\arraybackslash}p{0.41\linewidth}}
\toprule
ID & Planned Cell & Status & Evidence \\
\midrule
E1 & Independent + Field + Global & Implemented and run & \texttt{runs/phase2/E1\_Q*\_f*}, \texttt{scripts/run\_grid.py} \\
E2 & Independent + Field + RF-local & Implemented and run & \texttt{runs/phase2/E2\_Q*\_k*}, \texttt{scripts/run\_grid.py} \\
E3 & Independent + Vector + Global & Implemented and run & \texttt{runs/phase1/E3\_Q*\_f*}, \texttt{scripts/run\_grid.py} \\
E4 & Independent + Vector + RF-local & Implemented and run & \texttt{runs/phase2/E4\_Q*\_k*}, \texttt{scripts/run\_grid.py} \\
E5 & Incremental + Field + Global & Not fully implemented as planned grid cell & No E5 run directories; only partial Family II chain prototype on winner config \\
E6 & Incremental + Field + RF-local & Not implemented as planned grid cell & No E6 run directories \\
E7 & Incremental + Vector + Global & Not implemented as planned grid cell & No E7 run directories \\
E8 & Incremental + Vector + RF-local & Not implemented as planned grid cell & No E8 run directories \\
\bottomrule
\end{longtable}

\paragraph{Important nuance.}
A \emph{partial} Family II implementation exists in \texttt{src/train\_chain.py} and was executed in \texttt{runs/phase3b} and \texttt{runs/phase3b\_bptt}, but this is not the full E5-E8 factorial. It is winner-config chain-aware training on frozen Field SAEs.

\section{2) Thorough Summary of Implementation Changes in Actual Code}
The executed codebase diverges from the original all-8-cells-first intent into a staged, empirical narrowing strategy.

\subsection{A. Experiment flow changed from full factorial first to staged narrowing}
\begin{itemize}[leftmargin=*]
    \item Original plan: full E1-E8 grid.
    \item Actual flow:
    \begin{enumerate}[leftmargin=*]
        \item Phase 1 anchor (E3 only)
        \item Phase 2 branch (E1, E2, E4)
        \item Winner-driven hierarchy phases (Phase 3 and 3b)
    \end{enumerate}
    \item Operationalized in \texttt{scripts/run\_grid.py} and \texttt{scripts/run\_phase3.sh}.
\end{itemize}

\subsection{B. Family II implementation changed in form}
\begin{itemize}[leftmargin=*]
    \item High-level Family II describes incremental dependent sparse-code learning per adjacent pair with dependent code construction.
    \item Actual implementation: \texttt{src/train\_chain.py} performs chain-aware transition training with:
    \begin{itemize}[leftmargin=*]
        \item frozen winner SAEs,
        \item scheduled sampling,
        \item re-grounding (decode $\rightarrow$ re-encode $\rightarrow$ mask),
        \item optional BPTT through re-grounding.
    \end{itemize}
    \item This is a pragmatic Family II-style stabilizer, not full E5-E8 grid coverage.
\end{itemize}

\subsection{C. Added explicit causal upper-bound instrumentation (N4 hybrid)}
\begin{itemize}[leftmargin=*]
    \item \texttt{src/transitions.py} adds \texttt{frozen\_block\_hybrid}.
    \item This computes a causal upper bound by propagating sparse support through real frozen network blocks between sections.
    \item This mechanism is central to later conclusions (information sufficiency vs predictor instability).
\end{itemize}

\subsection{D. Added chain-mode diagnostics and recovery modes}
\begin{itemize}[leftmargin=*]
    \item \texttt{src/eval\_chain.py} evaluates:
    \begin{itemize}[leftmargin=*]
        \item raw learned chain,
        \item re-masked chain,
        \item re-grounded chain,
        \item frozen-block hybrid chain.
    \end{itemize}
    \item This directly exposes exposure-bias failure and the re-grounding mechanism.
\end{itemize}

\subsection{E. Extended beyond core E1-E8 with tier experiments}
Additional modules were implemented and run (at least partially):
\begin{itemize}[leftmargin=*]
    \item \texttt{src/intervene.py}: N8 parent-child causal intervention
    \item \texttt{src/feature\_audit.py}: feature health and duplicate/dead statistics
    \item \texttt{src/taxonomy.py}: sparse vs dense vs crosscoder transition taxonomy (implemented but run blocked)
    \item \texttt{src/vit\_sae.py}: ViT transfer pipeline (implemented but run blocked)
\end{itemize}

\subsection{F. Architecture and training details concretized}
\begin{itemize}[leftmargin=*]
    \item FieldSAE realized as ConvNeXt-style local residual blocks (5x5 depthwise, 1x1 mixing).
    \item VectorSAE realized as shared per-location 1x1 with decoder normalization.
    \item Hard TopK masking implemented for global and RF-local variants in \texttt{src/masks.py}.
    \item Layerwise training uses warmup $\rightarrow$ anneal $\rightarrow$ fixed curriculum in \texttt{src/train\_sae.py}.
\end{itemize}

\subsection{G. Practical engineering changes}
\begin{itemize}[leftmargin=*]
    \item Added environment workaround and orchestration scripts to handle CUDA/cuDNN and GPU health issues.
    \item Tier scripts show contingency workflows for dead GPU and node-state problems.
\end{itemize}

\section{3) Thorough Summary of Experiments Performed After Changes}
\subsection{Phase 1: E3 anchor reconstruction + controls}
\begin{itemize}[leftmargin=*]
    \item VectorSAE + Global mask + Independent protocol across Q1-Q5 and fractions.
    \item Random-dictionary controls run in \texttt{runs/phase1\_rand}.
    \item PCA baselines in \texttt{logs/baselines\_pca.json}.
\end{itemize}

\subsection{Phase 2: branch on architecture and mask}
\begin{itemize}[leftmargin=*]
    \item E1 (Field+Global), E2 (Field+RF-local), E4 (Vector+RF-local) across all 5 sections.
    \item Generated the practical winner configuration for hierarchy experiments.
\end{itemize}

\subsection{Phase 3: Family I transitions and chain evaluation}
\begin{itemize}[leftmargin=*]
    \item Adjacent transition models trained for Q1$\rightarrow$Q2, Q2$\rightarrow$Q3, Q3$\rightarrow$Q4, Q4$\rightarrow$Q5.
    \item Chain evaluated under raw, masked, re-grounded, and hybrid modes.
\end{itemize}

\subsection{Phase 3b: Family II-style chain-aware training}
\begin{itemize}[leftmargin=*]
    \item Scheduled-sampling training in \texttt{runs/phase3b}.
    \item BPTT-through-re-grounding variant in \texttt{runs/phase3b\_bptt}.
\end{itemize}

\subsection{Tier experiments (post-core)}
\begin{itemize}[leftmargin=*]
    \item Completed: N8 intervention, feature audits, Pareto sweep, stronger Q3$\rightarrow$Q4 stress test, ResNet-110 backbone training.
    \item Attempted but blocked by CUDA/node state: taxonomy run, ViT transfer run, R20/C10 recon grid, multi-seed grid, R110 recon grid.
\end{itemize}

\section{4) Results}
\subsection{Baseline backbone}
ResNet-56/CIFAR-100 frozen reference: \textbf{71.83\%} top-1 (from \texttt{logs/backbone\_r56\_c100\_\_done.json}).

\subsection{Core reconstruction findings (E1/E2/E3/E4)}
\subsubsection{Field vs Vector at aggressive global sparsity (2\% retained)}
Spliced top-1 from result JSONs:
\begin{itemize}[leftmargin=*]
    \item Q1: E1 0.7139 vs E3 0.6376
    \item Q2: E1 0.7164 vs E3 0.6327
    \item Q3: E1 0.7078 vs E3 0.6646
    \item Q4: E1 0.7106 vs E3 0.7098
    \item Q5: E1 0.7147 vs E3 0.7123
\end{itemize}

\subsubsection{RF-local viability (k\_rf=4) Field vs Vector}
\begin{itemize}[leftmargin=*]
    \item Q1: E2 0.7158 (frac 0.0312) vs E4 0.6858
    \item Q2: E2 0.6992 (frac 0.0078) vs E4 0.4645
    \item Q3: E2 0.6698 (frac 0.0039) vs E4 0.2430
    \item Q4: E2 0.6949 (frac 0.0078) vs E4 0.5942
    \item Q5: E2 0.6443 (frac 0.000122) vs E4 0.2504
\end{itemize}

\subsection{Phase 3 transition results}
Final values from \texttt{runs/phase3/T\_*/result.json}:
\begin{longtable}{>{\raggedright\arraybackslash}p{0.17\linewidth} >{\raggedright\arraybackslash}p{0.18\linewidth} >{\raggedright\arraybackslash}p{0.17\linewidth} >{\raggedright\arraybackslash}p{0.14\linewidth} >{\raggedright\arraybackslash}p{0.14\linewidth}}
\toprule
Transition & Learned spliced top-1 & Hybrid top-1 & Code MSE & relL2 \\
\midrule
Q1$\rightarrow$Q2 & 0.7146 & 0.7151 & 0.0078 & 0.2007 \\
Q2$\rightarrow$Q3 & 0.6907 & 0.7105 & 0.0437 & 0.4279 \\
Q3$\rightarrow$Q4 & 0.6504 & 0.7113 & 0.2005 & 0.5376 \\
Q4$\rightarrow$Q5 & 0.7089 & 0.7149 & 0.0586 & 0.4082 \\
\bottomrule
\end{longtable}

\subsection{Phase 3 chained results}
From \texttt{runs/phase3/chain\_result.json}:
\begin{itemize}[leftmargin=*]
    \item Raw learned chain: 0.0116
    \item Re-mask each step: 0.2126
    \item Re-ground each step: 0.5930
    \item Frozen-block hybrid chain: 0.7122
    \item Original reference top-1: 0.7183
\end{itemize}

\subsection{Phase 3b chain-aware results}
\begin{itemize}[leftmargin=*]
    \item Scheduled sampling (\texttt{runs/phase3b/result.json}):
    \begin{itemize}[leftmargin=*]
        \item raw chain: 0.0128
        \item re-grounded chain: 0.6624
        \item hybrid chain: 0.7122
    \end{itemize}
    \item BPTT re-grounding variant (\texttt{runs/phase3b\_bptt/result.json}):
    \begin{itemize}[leftmargin=*]
        \item raw chain: 0.0109
        \item re-grounded chain: 0.7010
        \item hybrid chain: 0.7122
    \end{itemize}
\end{itemize}

\subsection{Intervention and feature-quality results}
\subsubsection{N8 intervention (\texttt{runs/intervene\_result.json})}
\begin{itemize}[leftmargin=*]
    \item Importance ratio (top-ablate vs random-ablate):
    \begin{itemize}[leftmargin=*]
        \item Q1$\rightarrow$Q2: 3.29x
        \item Q2$\rightarrow$Q3: very large (random ablation near/below zero effect)
        \item Q3$\rightarrow$Q4: 4.00x
        \item Q4$\rightarrow$Q5: 7.50x
    \end{itemize}
    \item Locality fraction in RF neighborhood: about 0.38-0.50 vs null area fractions 0.009-0.141.
\end{itemize}

\subsubsection{Feature audit}
\begin{itemize}[leftmargin=*]
    \item Q3 (\texttt{runs/audit/audit\_Q3.json}): dead 2.73\%, near-dup 5.86\%.
    \item Q5 (\texttt{runs/audit/audit\_Q5.json}): dead 1.37\%, near-dup 0.00\%.
\end{itemize}

\subsection{Baselines}
\begin{itemize}[leftmargin=*]
    \item Random dictionary controls (\texttt{runs/phase1\_rand}): near chance on deep sections (for example Q3 0.0093, Q4 0.0097, Q5 0.0116).
    \item PCA baseline (\texttt{logs/baselines\_pca.json}) substantially below trained SAE at comparable or larger budgets in many sections.
\end{itemize}

\section{Blocked/Pending Runs and Their Status}
Based on \texttt{runs/tier\_rest/*.log} and grid summaries:
\begin{itemize}[leftmargin=*]
    \item \texttt{runs/r20c10/grid\_summary.json}: 15/15 failed (no result.json outputs).
    \item \texttt{runs/seeds/grid\_summary.json}: 15/15 failed (no result.json outputs).
    \item \texttt{runs/tier\_rest/taxonomy.log}: CUDA initialization unavailable.
    \item \texttt{runs/tier\_rest/vit\_sae.log}: CUDA initialization unavailable.
    \item \texttt{runs/tier\_rest/r110\_Q1.log}: failed due missing cached acts file (dependency chain interrupted after cache failure).
\end{itemize}

This is consistent with the documented node/GPU fault context in README/TODO/report notes.

\section{Final Answers to the Four Requested Questions}
\subsection*{Q1: Which of 8 are implemented?}
Implemented and run: \textbf{E1, E2, E3, E4}.\newline
Not implemented as full planned cells: \textbf{E5, E6, E7, E8}.

\subsection*{Q2: What changed in implementation?}
The project moved from a planned full 8-cell factorial to a staged winner-first pipeline; implemented robust chain diagnostics and a causal upper bound; implemented Family II as chain-aware scheduled-sampling/re-grounding on winner SAEs rather than full E5-E8 factorial training.

\subsection*{Q3: What experiments were performed after changes?}
Phase 1 (E3 + controls), Phase 2 (E1/E2/E4), Phase 3 transitions and chained evaluations, Phase 3b chain-aware/BPTT improvements, plus tier experiments (intervention, feature audit, Pareto, stress tests), with some higher-tier runs blocked by runtime GPU/node issues.

\subsection*{Q4: Results}
Core numerical outcomes are:
\begin{itemize}[leftmargin=*]
    \item Strong sparse reconstruction/preservation with Field SAE at low retained fractions.
    \item Transition single-step performance high except pronounced Q3$\rightarrow$Q4 gap.
    \item Raw learned 4-step chain collapses to chance, but re-grounding recovers major performance; BPTT re-grounding reaches 0.701 vs 0.712 hybrid bound.
    \item Intervention and feature-audit outcomes support non-trivial, structured sparse representations.
\end{itemize}

\section{Appendix: Key Paths for Re-Verification}
\begin{itemize}[leftmargin=*]
    \item Planned grid definitions: \texttt{high-level/sparse\_signal\_sae\_experiment\_plan\_high\_accuracy\_conv.tex}
    \item Actual job definitions: \texttt{scripts/run\_grid.py}
    \item Transition training: \texttt{src/train\_transition.py}
    \item Chain eval and variants: \texttt{src/eval\_chain.py}
    \item Family II-style training: \texttt{src/train\_chain.py}
    \item E1-E4 artifacts: \texttt{runs/phase1}, \texttt{runs/phase2}
    \item Phase 3 artifacts: \texttt{runs/phase3}
    \item Phase 3b artifacts: \texttt{runs/phase3b}, \texttt{runs/phase3b\_bptt}
    \item Intervention and audit artifacts: \texttt{runs/intervene\_result.json}, \texttt{runs/audit/*}
\end{itemize}

\end{document}
